%! AP Statistics — Unit 4, Lesson 4.1 — Group Activity KEY
\documentclass[10pt]{article}
\usepackage{apstats-article}
\usepackage{apstats-key}

\begin{document}

\pageheader{Unit 4, Lesson 4.1}{Group Activity --- KEY}
\namepartnerperiod

\vspace{0.1in}

\begin{scenariobox}[Context]{sky}
\small
For each scenario below, your group will (1) identify the pattern observed, (2) write a
statistical question it suggests, (3) decide whether chance alone is a plausible
explanation, and (4) describe what additional evidence would confirm a non-random cause.
Work through all tiers your teacher assigns.
\end{scenariobox}

\vspace{0.1in}

% ── Tier R ────────────────────────────────────────────────────────────────────
\begin{tcolorbox}[colback=goldbg, colframe=goldacc, boxrule=1pt, arc=2mm,
    left=4mm, right=4mm, top=3mm, bottom=3mm,
    title={\bfseries Tier R \enspace Remediate --- Getting Started with Patterns}]
\small

\textbf{Scenario 1:} A student rolls a standard die 20 times and records:
\[
1,\ 3,\ 5,\ 5,\ 5,\ 2,\ 4,\ 6,\ 5,\ 1,\ 3,\ 2,\ 5,\ 5,\ 4,\ 6,\ 1,\ 2,\ 3,\ 5
\]

\begin{enumerate}[leftmargin=1.5em, itemsep=8pt]
  \item \ans{5 appeared 7 times.}
  \item \ans{More than expected; a fair die should show 5 about $20 \times \frac{1}{6} \approx 3.3$ times. Getting 7 is higher than average.}
  \item \ans{This pattern \textbf{suggests a question} about the die. (Pattern alone does not prove unfairness.)}
  \item \ans{Is the proportion of 5s rolled by this die significantly higher than $\frac{1}{6}$, or could this result occur by chance with a fair die?}
\end{enumerate}

\vspace{8pt}
\textbf{Scenario 2:} Your school basketball team wins 8 games in a row after losing
their first 3 games of the season.

\begin{enumerate}[leftmargin=1.5em, itemsep=8pt]
  \item \ans{The team lost 3 in a row, then won 8 consecutive games.}
  \item \ans{No --- streaks of 8 wins can happen by chance, especially if the team is already playing against weaker opponents.}
  \item \ans{Is the team's win probability different in the second part of the season compared to the first, or could an 8-game winning streak occur by chance?}
\end{enumerate}

\end{tcolorbox}

\vspace{0.1in}

% ── Tier A ────────────────────────────────────────────────────────────────────
\begin{tcolorbox}[colback=greenbg, colframe=greenacc, boxrule=1pt, arc=2mm,
    left=4mm, right=4mm, top=3mm, bottom=3mm,
    title={\bfseries Tier A \enspace Approaching Proficiency --- Patterns and Statistical Questions}]
\small

\textbf{Scenario 3:}

\begin{enumerate}[leftmargin=1.5em, itemsep=8pt]
  \item \ans{$11/400 = 0.0275 = 2.75\%$}
  \item \ans{The observed rate (2.75\%) is about 5.5 times the national rate (0.5\%).}
  \item \ans{Is the rate of this liver condition in this neighborhood significantly higher than the national rate, or could the observed cluster be due to random variation?}
  \item \ans{Possibly --- small populations can produce apparent clusters by chance. With an expected count of only 2 cases (400 × 0.5\%), getting 11 is quite unlikely by chance and deserves investigation.}
  \item \ans{Multiple years of data showing elevated rates; identification of a shared environmental exposure; comparison with neighboring areas; a controlled epidemiological study.}
\end{enumerate}

\vspace{8pt}
\textbf{Scenario 4:}

\begin{enumerate}[leftmargin=1.5em, itemsep=8pt]
  \item \ans{Likes per hour jumped from about 12,000 to 80,000 (nearly 7×) starting at hour 7.}
  \item \ans{Is the like rate during hours 7--8 significantly higher than the rate during hours 1--6, or could this spike be due to random fluctuation in social media engagement?}
  \item \ans{No --- a nearly 7× jump in rate is very unlikely to occur by chance in a stable process. A more plausible explanation is a retweet by another large account, a trending hashtag, or a media mention.}
\end{enumerate}

\end{tcolorbox}

\vspace{0.1in}

% ── Tier E ────────────────────────────────────────────────────────────────────
\begin{tcolorbox}[colback=sky, colframe=navy, boxrule=1pt, arc=2mm,
    left=4mm, right=4mm, top=3mm, bottom=3mm,
    title={\bfseries Tier E \enspace Extension --- Deeper Analysis}]
\small

\textbf{Scenario 5:}

\begin{enumerate}[leftmargin=1.5em, itemsep=8pt]
  \item \ans{9 consecutive days above 90°F.}
  \item \ans{Is the probability of observing 9 or more consecutive days above 90°F in a 30-day period significantly higher than what historical temperature patterns would predict, or could this streak be explained by natural random variation?}
  \item \ans{If temperatures are positively autocorrelated (not independent), then the probability of a long run of hot days is higher than a simple independent model predicts. This means the threshold for calling the streak ``non-random'' needs to account for the dependence structure. Random does not mean independent --- the skeptic raises a valid statistical concern.}
  \item \ans{Compare the frequency of 9-day (or longer) runs above 90°F in historical records from before a suspected climate change period vs.\ after. Use many years of data. A formal test would compare observed run frequencies to what a model of natural temperature variability would predict.}
\end{enumerate}

\end{tcolorbox}

\end{document}
